\chapter{Trabalhos Relacionados}
\label{cap:trabalhos-relacionados}

A literatura sobre análise automatizada de textos para diagnóstico de políticas públicas e avaliação de sustentabilidade tem crescido significativamente nos últimos anos, impulsionada pela disponibilidade de grandes volumes de dados não estruturados e pelo avanço dos Modelos de Linguagem de Grande Porte (\gls{LLM}s). Esta seção apresenta os principais trabalhos que contribuem para o embasamento teórico e metodológico desta dissertação, organizados em quatro eixos temáticos: análise de sentimentos aplicada a dados públicos, uso de \gls{NLP} para avaliação de políticas de sustentabilidade, aplicações de LLMs em contextos urbanos e o uso de dados do Google Places como fonte de diagnóstico.

\section{Análise de Sentimentos Aplicada a ESG e Políticas Públicas}
\label{sec:tr-esg}

\citeauthor{liu2023esg} \citeyear{liu2023esg} investigaram percepções públicas sobre fatores ambientais, sociais e de governança (\gls{ESG}) a partir de dados de mídia social coletados na China. O estudo utilizou técnicas de análise de sentimentos sobre textos de plataformas digitais e demonstrou que a opinião pública capturada de forma não estruturada pode antecipar e complementar indicadores ESG formais. A metodologia bem estruturada e os resultados validados demonstraram que a análise de sentimentos pode revelar percepções públicas que dificilmente seriam capturadas por métodos tradicionais de pesquisa. Este trabalho se relaciona diretamente com a presente dissertação ao validar o uso de dados textuais públicos para inferência sobre comportamentos coletivos ligados ao desenvolvimento sustentável.

\citeauthor{schimanski2024} \citeyear{schimanski2024} propuseram um framework para quantificar comunicações ESG utilizando \gls{NLP}. Os autores utilizaram modelos baseados em Transformers, como \gls{RoBERTa} e DistilRoBERTa, para analisar mais de 13,8 milhões de textos corporativos, demonstrando que técnicas de NLP podem medir com confiabilidade o nível de conformidade e transparência em reportes de sustentabilidade. A contribuição central deste trabalho é a evidência empírica de que modelos pré-treinados de linguagem podem ser adaptados para domínios altamente especializados, como a comunicação ESG, sem necessidade de \textit{fine-tuning} extensivo.

\citeauthor{jaiswal2024} \citeyear{jaiswal2024} analisaram o sentimento do mercado financeiro em relação a investimentos ESG por meio de mineração de opinião em dados do Twitter. O estudo revelou padrões temáticos e emocionais relevantes para a compreensão de como o público percebe iniciativas de sustentabilidade corporativa. Apesar de focar no contexto financeiro, o estudo contribui metodologicamente para a presente pesquisa ao demonstrar a viabilidade de análise de sentimentos em larga escala sobre textos curtos e heterogêneos.

\section{Uso de NLP para Avaliação de Sustentabilidade e ODS}
\label{sec:tr-ods-nlp}

\citeauthor{gupta2024} \citeyear{gupta2024} desenvolveram um modelo de processamento de linguagem natural baseado em \gls{BERT} combinado com a técnica YAKE para extração de palavras-chave em relatórios de sustentabilidade corporativa. O modelo alcançou 98\% de precisão na extração de termos relevantes, confirmando a eficácia de abordagens híbridas que combinam modelos neurais profundos com técnicas estatísticas de extração. O trabalho é relevante para a presente dissertação por demonstrar como modelos de NLP podem ser especializados para extrair informações estruturadas de textos sobre sustentabilidade.

\citeauthor{larosa2025} \citeyear{larosa2025} investigaram o potencial e as limitações dos LLMs na análise de políticas climáticas e de sustentabilidade. Os autores argumentam que LLMs genéricos oferecem vantagens significativas em contextos de baixa disponibilidade de dados rotulados, podendo realizar análise temática e identificação de lacunas em documentos de política pública sem necessidade de treinamento específico. Esta perspectiva é central para a justificativa metodológica da presente dissertação, que adota modelos de linguagem generativos (Gemini) para análise de textos públicos sobre ODS.

O trabalho de \citeauthor{abidoye2022} \citeyear{abidoye2022} discute a localização dos ODS por meio de dados desagregados, argumentando que indicadores nacionais ou estaduais frequentemente obscurecem disparidades locais críticas. Os autores defendem o investimento em coleta de dados granulares a nível municipal para avaliação mais precisa do progresso em direção às metas da Agenda 2030. Esta perspectiva reforça a relevância metodológica desta dissertação, que trata Mossoró-RN como unidade de análise e utiliza dados georreferenciados coletados no nível do estabelecimento.

\section{LLMs e Dados Urbanos para Diagnóstico de Políticas}
\label{sec:tr-llm-urbano}

\citeauthor{fu2023} \citeyear{fu2023} propuseram um framework de ciência urbana colaborativa baseado em LLMs pré-treinados, demonstrando que modelos como \gls{GPT} podem ser utilizados para análise de fenômenos urbanos complexos, incluindo mobilidade, habitação e uso do solo. O estudo demonstrou a viabilidade de integrar dados heterogêneos de fontes urbanas com análise gerada por IA, produzindo insights acionáveis para planejamento urbano. A abordagem proposta serve de referência para o uso de Gemini na presente dissertação para síntese e diagnóstico baseados em avaliações de usuários sobre estabelecimentos em Mossoró-RN.

\citeauthor{wankhade2022} \citeyear{wankhade2022} realizaram uma revisão abrangente de métodos de análise de sentimentos, aplicações e desafios. Os autores categorizaram as abordagens em três grupos: baseadas em léxico, baseadas em aprendizado de máquina tradicional e baseadas em aprendizado profundo. A revisão identificou que modelos baseados em Transformers (BERT, RoBERTa) representam o estado da arte em precisão, enquanto LLMs generativos emergem como alternativa flexível para contextos onde dados rotulados são escassos. Esta taxonomia orienta as escolhas metodológicas desta dissertação.

\citeauthor{agua2025} \citeyear{agua2025} investigaram o uso de LLMs para análise de sentimentos baseada em aspectos (\textit{aspect-based sentiment analysis}) em avaliações de clientes no setor de turismo. O estudo demonstrou que modelos de linguagem generativos podem identificar, sem treinamento adicional, sentimentos associados a aspectos específicos mencionados em textos livres -- exatamente a abordagem empregada nesta dissertação para classificar sentimentos relacionados às dimensões dos ODS 9 e ODS 16.

\section{Google Places como Fonte de Dados para Análise Urbana}
\label{sec:tr-google-places}

A utilização da Google Places API como fonte de dados para análise urbana e de políticas públicas é uma abordagem emergente que combina a ampla cobertura geográfica da plataforma com a riqueza qualitativa das avaliações de usuários. Diferentemente de pesquisas domiciliares ou questionários estruturados, as avaliações no Google Places refletem experiências espontâneas e não mediadas de cidadãos com estabelecimentos públicos e privados, tornando-as uma fonte valiosa para capturar percepções sobre qualidade de serviços, infraestrutura e segurança pública.

O IDSC -- Índice de Desenvolvimento Sustentável das Cidades \cite{idsc2024} -- calcula anualmente pontuações para os 17 ODS em municípios brasileiros a partir de indicadores primários como dados do IBGE, registros administrativos e pesquisas oficiais. Embora robusto para análises comparativas nacionais, o IDSC não fornece informações granulares sobre as causas específicas do baixo desempenho em determinado ODS. Esta dissertação propõe uma abordagem complementar ao IDSC, utilizando avaliações públicas do Google Places para diagnosticar, com especificidade local, os fatores que explicam as pontuações observadas.

Em síntese, os trabalhos revisados confirmam a viabilidade e relevância das escolhas metodológicas desta dissertação: a combinação de coleta automatizada de avaliações públicas, filtragem semântica multinível e análise por LLMs para diagnóstico de ODS é inovadora, replicável e alinhada com o estado da arte em análise de sentimentos, NLP para sustentabilidade e ciência urbana baseada em dados.
